<<P000640>
<extensionDeclaration> ::= <<<P001066><
<<<P001067>{} <typeIdentifierNotType> <typeParameters> \ON{{<тип>>
\gnewline{}\{" (<metadata> <classMemberDeclaration>)* <<<P001527> "
<<P000658>

<<<P000958>%
Декларация, полученная из <<<P001070>{продление Декларации}, известна как
<<<P001071>{продление!декларация}.
Он вводит в действие
<<P000747>
с данным <<<P001072>{идентификатор}, если присутствует,
в пространство имен библиотеки, в которой находится
<<<P001073>{(и, следовательно, в библиотечном объеме)},
и предоставляет возможность для объявления
членов-расширителей.
Дополнительно расширение вводится в пространство имен библиотеки
Свежее, личное имя.
Первый из них известен как
<<<P001074>{объявленное имя}{расширение!объявленное имя}
от продления,
Последний известен как
<<<P001075>{свежее имя}{расширение! свежее название
расширения.

<<<P001076>{%>
Также вводится новое название.
Пространство имен библиотеки текущей библиотеки
для каждого импортного расширения,
Даже если он импортируется с приставкой
(P000675).
%
Заявленное название расширения <<<P000016> это
введены в библиотечный объем текущей библиотеки
По тем же правилам, что и названия
Другие местные или импортные декларации, такие как классы.
%
Оригинальное название: <<P000017> вводятся в этих случаях,
Но и в одном дополнительном случае:
когда происходит столкновение имени на объявленном имени <<P000018>.%
?

<<<P001077>{%>
Новое название позволяет использовать расширение
в неявном призыве
(<<P000676>),
Даже в том случае, если заявленное имя
или префикс импорта, обеспечивающий доступ к заявленному названию
оговаривается декларацией в промежуточном объеме,
Конфликт с именем, столкновение.%
?

<<<P001078>{%>
Тот факт, что расширение может быть использовано неявно даже в том случае, когда
не имеет объявленного имени
или объявленное имя затенено или противоречиво
отражает тот факт, что основное предполагаемое использование является неявным вызовом.
Хотя разработчик не может знать (и, следовательно, не может использовать) новое имя.
из данного расширения имплицитное обращение может использовать его.%
?

<<<P000959>%
Ошибка времени компиляции, если библиотека имеет
Отсроченный импорт библиотеки <<<P000019>
Импортируемое пространство имен из <<<P000020> содержит
Имя, обозначающее расширение.

<<<P001079>{%>
Это означает, что импорт должен \HIDE{} или <<P001081>>>{} Для устранения
Наименования расширений от отложенного импорта.%
?

<<<P001082>{%>
Это ограничение гарантирует, что никакие расширения не будут введены.
использование отсроченного импорта,
что позволяет ввести семантику для таких расширений в будущем.
Без ущерба для существующего кода.%
?

<<<P000960>%
<<<P001083>{тип} в заявлении о продлении
называется расширением
\IndexCustom{\ON{} Тип {расширение!\ON{} Тип.
<<<P001087> Тип может быть любым действительным типом, включая переменную типа.

<<<P001088>{%>
Основная интуиция заключается в том, что расширение <<P000767> может иметь
Тип <<<P001089>{} <<P000021> (определение типа получателя) и набора членов.
Если <<P000022> выражение, статический тип которого <<<P000023>
<<P000768> является членом, объявленным <<<P000769>,
<<P000770> может ссылаться на указанный член со значением <<<P000025> Ссылка на <<<P001090>.
Явно разрешенная форма <<<P000771> Доступен,
<<<P000772> Его можно использовать даже в случае
где <<P000773> Выполняет какую-то другую функцию
Потому что некоторые другие расширения более специфичны.
Приведены подробности этих понятий, правил и механизмов.
в этом разделе и его подразделах.%
?

<<P000961>%
Заявленное название расширения не обозначает тип,
Можно использовать для обозначения самого расширения.
<<<P001091>{%>
(например, для доступа к статическим элементам расширения,или для того, чтобы разрешить вызов явно)%
}

<<P000962>%
Декларация о продлении вводит две сферы:

<<P000641>
<<P001092>
A <<<P001093>{тип-параметр}{скоп!тип-параметр},
который является пустым, если расширение не является общим (<<P000677>).
Сфера охвата заявления о продлении по типу параметра включает
библиотечный объем текущей библиотеки.
Сфера действия параметра типа - это текущая сфера для
Параметры типа и для расширения <<<P001094> <<<P001095>{тип}.
<<P001096>
A <<<P001097>{body scope}{scope!extension body}.
Включающая сфера охвата органа сфера действия заявления о продлении является
Сфера действия декларации о продлении типа.
В настоящее время сфера действия заявления о продлении членства является
сфера охвата прилагаемой декларации о продлении.
<<P000659>

<<P000963>%
<<<P001047>%
Пусть <<<P000028> будет заявлением о продлении с объявленным именем <<<P000775>.
Декларация в <<<P000029> с модификатором <<<P001098>{} обозначается как
<<<P001099>{статический член}{расширение!статический член}
Декларация.
%
Декларация в <<<P000030> без модификатора \STATIC{} обозначается как
<<<<P001101>{член организации}{расширение!член организации}
Декларация.
%
Ошибки в названии конфликтов могут возникать в <<<P000031>
В ситуациях, которые также происходят в классах и миксах
(P000678).
Кроме того, ошибка времени компиляции происходит в следующих ситуациях:

<<P000642>
<<<p001102> <<P000032> Объявляет участника, имя основания которого <<P000776>.
<<P001103> <<P000033> Заявляет параметр типа <<<P000777>.
<<P001104> <<P000034> Объявляет участника, базовым именем которого является имя параметра типа
<<<P000035>.
<<P001105> <<P000036> объявляется экземпляром члена или статичным членом, имя основания которого
<<P000778>, <<P000779>, <<P000780>, <<P000781>,
или <<<P001106>{==}.
<<<P001107>{%>
То есть член, чье имя также является именем
Член экземпляра, который есть у каждого объекта.%
?
<<P001108> <<P000037> объявляется застройщиком.
<<P001109> <<P000038> объявляет переменную экземпляра.
<<P001110> <<P000039> Объявляют абстрактным членом.
<<P001111> <<<P000040> Объявляет метод с формальным параметром
с модификатором <<<P001112>.
<<P000660>

<<<P001113>{%>
Абстрактные члены не допускаются, поскольку
Никакой поддержки для реализации.
%
Конструкторы не допускаются с момента продления
Не вводит никаких типов, которые могут быть построены.
%
Переменные не допускаются, поскольку память не выделяется.
для каждого объекта, доступного как <<<P001114>{} в членах.
Разработчики могут эмулировать per-\THIS государство при необходимости,
Например, используя <<<P000782>.
%
Члены с тем же именем базы, что и члены <<<P000783>
не допускаются, поскольку они могут быть вызваны только с использованием явного разрешения;
как в <<P000784>,
Что может быть запутанным и подверженным ошибкам.%
?


<<<P001116> Явное прошение члена представительства о продлении
<<P001032>

<<P000964>%
<<<P001117> E - простой или квалифицированный идентификатор
Это означает расширение.
<<<P001118>{приложение расширения}{приложение расширения!
является выражением формы
<<<P001119><<<P000041> <typeArguments>? '(' <expression> ')'.
Член расширения с именем, доступным для текущей библиотеки
может быть явно использован на конкретном объекте
Выполняя призыв члена
(<<P000679>)
где приемник является приложением расширения.

<<<P001120>{%>
Тип вывода еще не указан в этом документе,
Предполагается, что они уже имели место
(<<P000680>),
Ниже приводится описание предполагаемого лечения.
В этом разделе и его подразделах есть похожие комментарии о выводе типа.
Ниже, с пометкой «С выводом типа: <<<P001121>».

<<<P000042> быть простым или квалифицированным идентификаторомРасширение под названием <<P000785> и заявлены следующим образом:%
?

<<P000000>

<<P001122>
<<<P001123>{%>
Тип вывода для приложения расширения
в форме <<P000786>
делается точно так же, как это было бы для
Тот же синтаксис, рассматриваемый как вызов конструктора
где <<<P000045> предполагается обозначать следующий класс,
и тип контекста пуст (без требований): %
?

<<<P000001>

<<<P001124>{%>
Это выведет аргументы типа для <<<P000787>, и это будет
Введите тип контекста для выражения <<P000048>.
Например, если <<<P000049> объявляется как
\code{<<P001126>>>{) E<T> в Set<T> \{\,\ldots\,\}
<<<P001128> E(\{\})} будет содержать выражение \code{\{\}} с
тип контекста, который делает его буквальным.%
?

<<P000965>%
<<<P001048>%
<<<P000050> быть расширением формы
<<<P001130><<<P000051><<<<P001131> T}{1}{s}>($e$],
где <<P000053> обозначает расширение, объявленное как

<<P001132>
<<<P001049>%
\EXTENSION{} E<<<<P001134> <<<P001135>< <<P000054> \{\,\ldots\,<<<P001539>.

<<P000966>%
Мы определяем
<<<P001137><обоснованно \ON{} {расширение!обосновано} \ON{} тип
<<<P000055> <<<P000056>.
Мы определяем
<<<P001140>{обоснование-к-связанности <<<P001141>< Тип: {%}
Обсуждение:Instantiation-to-bound <<<P001142>< тип
<<<P000057> <<<P000058>
где <<<P001143>{U}{1}{s} является результатом инстанциации в связанную
Типовые параметры <<<P000059>
(P000681).
%
Ошибка времени компиляции возникает, если
<<P001144> T_j}{[T_1/X_1, \ldots, T_s/X_s]B_j},
<<<p000060>
<<<P001146>{(то есть границы не могут быть нарушены)}.
%
Ошибка времени компиляции возникает, если статический тип <<<P000061> может быть присвоено
<<<P001147>> Тип <<<P000062>.
<<<P001148>{%>
Обратите внимание, что ошибка времени компиляции также происходит.
если статический тип <<<P000063> является <<<P001149>
(<<P000682>)%
?

<<<P000967>%
Это ошибка времени компиляции, если происходит расширение приложения.
в месте, где это <<<P001150>{не} синтаксический приемник
простой или составной член вызов
(P000683).

<<<P001151>{%>
То есть, единственным допустимым использованием приложения расширения является
Призывать или отрывать от него членов.
Это похоже на то, как имена префиксов могут использоваться только в качестве
Цели призыва членов,
Кроме того, расширения могут декларировать операторы.
Например, <<P000788> может быть действительным призывом к
Оператор <<<P001152>{+} объявлено в расширении <<<P000066>.%
?

<<<P000968>%
Приложение расширения не имеет типа.

<<<P001153>{%>
Это согласуется с тем, что любое использование приложения расширения
где тип необходим - ошибка времени компиляции.
?

<<P000969>%
<<<P001050>%
<<<P000067> Быть простым членом призыва
(<<P000684>)
чей приемник <<P000068> является расширением формы
<<<P001051>%
<<<P001154><<<P000069><<<<P001155> T}{1}{k}>($e$)}
\commentary{(который равен <<<P000789>, когда <<<P000073> равен нулю)}
имя члена-корреспондента <<<P001157>{n},
Предположим, что <<P000074> Не имеет ошибок времени компиляции.
Ошибка времени компиляции возникает, если расширение не обозначено <<<P000075>
Объявляют члена по имени <<P000076>.
В противном случае пусть <<<P001158><<<P001159> X}{1}{k}}
Типовые параметры указанного расширения.
Пусть \DefineSymbol{s} будет подписью члена <<<P000077> Об этом сообщает <<<P000078>.
Точно такие же ошибки времени компиляции происходят для <<<P000079> как
те, которые произошли бы для вызова члена <<<P000080>
который получен из <<<P000081> Заменить <<P000082> через
переменной, тип которойКласс <<<P000083> в том же объеме, что и <<P000084>
что объявляет члена <<<P000085> с подписью члена
<<P000086>:

<<P000012>

<<<P000970>%
Подпись участника <<<P000087> называется
\IndexCustom{подпись участника вызова}{%}
Расширение! подпись участника вызова.
<<P000088>.
Статический тип <<<P000089> является возвратным типом
подпись участника вызова <<<P000090>.
<<P001162> Например: }

<<P000002>

<<<P001163>{%>
С выводом типа:
В случае, когда подпись участника вызова <<<P000091> является родовым,
Вывод типа происходит на <<<P000092> Точно так же, как это
Для вызова функции, тип которой является
Тип функции <<<P000093>.%
?

<<<P000971>%
<<<P001052>%
Для динамической семантики,
<<<P000094> Быть простым, безусловным членом
чей приемник <<P000095> является расширением формы
<<<P001053>%
<<<P001164><<<P000096><<<<P001165> T}{1}{k}> (<<<P000097>)
где параметры типа <<<P000098> \DefineSymbol{\List{ X}{1}{k}}
и фактические значения \List{T}{1}{k} являются <<P001169>>>{<<P001170>>>{t}{1}{k}}
(<<P000685>),
<<<P001054>%
и имя члена-корреспондента которого <<<P000099>.
<<<P000100> быть членом <<<P000101> Называется он «P000102».
Оценка <<<P000103> доходов путем оценки
<<P000104> для объекта <<<P000105>,
оценка и привязка фактических аргументов к формальным параметрам;
(<<P000686>),
и, наконец, выполнить <<<P000106>
в среде связывания, где \List{X}{1}{k} связаны с \List{t}{1}{k},
<<<P001173>{} связано с <<<P000107>,
Каждый формальный параметр связан с соответствующим фактическим аргументом.
Значение <<<P000108> Это значение, возвращаемое вызовом <<<P000109>.

<<P000972>%
Когда <<P000110> является условным или составным призывом,
Статический анализ и динамическая семантика определяются
член призыва разбавлять
(P000687).

<<<P001174>{%>
Обратите внимание, что каскад (<<P000688>)
Приемник приложения расширения <<<P000111> представляет собой ошибку времени компиляции.
Это происходит потому, что подразумевается, что <<P000112> обозначает объект, который не является истинным,
а также потому, что это заставит каждого <<<P001175>
для вызова члена того же расширения,
который вряд ли будет желательным.%
?


<<<P001176> Неявное прошение члена представительства о продлении
<<P001033>

<<P000973>%
Вспомогательные элементы расширения могут быть вызваны или скрытно закупорены.
<<<P001177>{(без упоминания названия расширения)},
Как если бы они были членами получателя
Призыв данного члена.

<<<P001178>{%>
Например, если <<<P000790> это
правильное явное обращение члена экземпляра <<<P000791> из
расширение <<<P000792>,
<<<P000793> может быть
Правильное неявное обращение с тем же значением.
Другими словами,
Получатель <<<P000120> вызова участника может быть неявно заменен
приложение расширения с приемником <<<P000121>,
если удовлетворяется ряд требований, описанных ниже.%
?

<<<P001179>{%>
Неявное вызов предназначен в качестве основного способа использования расширений.
с явным призывом в качестве запасного варианта
если неявное обращение является ошибкой, или
Неявное обращение решается в отношении члена
расширение, отличное от предполагаемого.%
?

<<P000974>%
<<<P001055>%
Происходит неявное вызов члена-расширителя
Призыв к члену <<P000122>
(перенаправлено с «P000689»)
с приемником <<P000123> имя соответствующего члена <<<P000124> если {f}
(1) <<P000125> не является буквальным,
(2) интерфейс статического типа <<<P000126> не имеет члена
имя основания которого является именем основания <<<P000127>,
3 существует уникальный, наиболее специфический
(<<P000690>)Расширение, обозначаемое <<<P000128> который является доступным
(<<P000691>)
и применимым
(<<P000692>)
<<P000129>.

<<P000975>%
В случае, если ошибка компиляции не произошла,
<<P000130> Обозначается как <<<P000131>, который получен из <<<P000132> через
Замена ведущего <<<P000133> на <<<P000794>.

<<<P001180>{%>
С выводом типа:
Когда <<P000136> является родовым,
Типовой вывод, применяемый к <<<P000795> могут приводить реальные аргументы,
Получение <<<P000139> форма \code{$E$<\List{ T}{1}{k}>(<<P000141>>>)}.
Если этот этап вывода типа не удается, то <<<P000142> не применяется
(<<P000693>) %
?

<<<P001183>{%>
Неявное обращение члена инстанции о продлении
В каскаде также возможно
потому что каскад разбавляется до выражения, содержащего
одно или несколько обращений.%
?


<<<P001184> Доступность расширения.
<<<P001034>

<<P000976>%
Расширение <<<P000143> это
<<<P001185>{доступно}{расширение!доступность}
в определенной сфере <<P000144>
Если существует имя <<P000145> Для лексического поиска <<<P000146> из <<<P000147>
(<<P000694>)
доходность <<P000148>.

<<<P001186>{%>
Оригинальное название: P000149 может быть объявлено имя <<<P000150> или свежее название <<<P000151>,
Но так как свежее название всегда в объеме
Всякий раз, когда заявленное имя находится в пределах,
Достаточно вспомнить свежее название.
%
Когда новое имя <<P000152> находится в библиотечном пространстве,
доступно в формате <<<P001187>{любой},
потому что имя свежее и потому не может быть затенено
любым заявлением в любой промежуточной сфере.
%
Это означает, что если <<P000153> Доступен в любом месте данной библиотеки <<<P000154>
Тогда он доступен везде в <<<P000155>.%
?


<<P001188> Применимость продления
<<<P001035>

<<P000977>%
<<<P001056>%
Пусть <<P000156> будет расширением.
<<<P000157> Быть членом призыва
(<<P000695>)
с приемником <<<P000158> Статический тип <<<P000159>
и с именем соответствующего члена, базовое имя которого <<<P000160>.
Мы говорим, что <<P000161> это
<<<P001189>{применимо}{расширение!применимо}
<<<P000162>> Если соблюдены все следующие условия:

<<P000643>
<<P001190>
<<P000163> Это не буквальный тип.

<<<P001191>{%>
Это означает, что призыв является вызовом члена экземпляра,
В том смысле, что <<<P000164> обозначает объект,
Таким образом, он может вызвать элемент экземпляра или быть ошибкой,
Это не может быть статический доступ.
Обратите внимание, что <<<P000165> Также не обозначает префикс или расширение,
И это не расширение,
Потому что у них нет типа.%
?
<<P001192>
Тип <<<P000166> не имеет экземпляра с именем базы <<<P000167>,
<<<P000168> не является ни <<<P001193>{}, ни <<<P000796>.

<<<P001194>{%>
\DYNAMIC{{} и <<P000797>>> Считается, что у них есть все члены.
Кроме того, это ошибка для доступа к члену на приемнике типа. <<<P001196><
(<<P000696>),
Расширения никогда не применяются к получателям
любой из типов <<<P001197>, <<P000798> или <<P001198>.%
?

Для определения применимости расширения,
Типы функций и тип <<P000799>
Считается, что у них есть член по имени <<<P001199>.

<<<<P001200><%
Следовательно, расширения никогда не применимы к функциям.
когда базовое имя члена <<<P001201>.
%
Участников, заявленных встроенным классом <<<P000800>
существуют на всех видах,
Для членов с такими именами расширение никогда не применяется.%
?
<<<p001202>
Рассмотрим приложение расширения <<<P000169> формы <<<P000801>,
где <<P000172> Новая переменная со статическим типом <<P000173>.
Необходимо, чтобы возникновение <<<P000174>в области, которая является текущей областью для <<<P000175>
Это не ошибка времени компиляции.

<<<P001203>{%>
Другими словами, <<<P000176> должен соответствовать <<<P001204> Тип <<<P000177>.
С выводом типа могут быть добавлены выведенные фактические аргументы типа,
доходность \code{<<P000178>>><\List S}{1}{k}>($v$],
Который не должен быть ошибкой.
Если этот шаг вывода не срабатывает, это не ошибка,
Это означает, что <<P000180> Не применяется к <<<P000181>.%
?
<<P001207>
Расширение <<<P000182> Заявляет экземпляр члена с именем базы <<<P000183>.
<<P000661>

<<<P001208>{%>
С выводом типа:
Контекстный тип обращения не влияет
применимо ли продление,
и ни тип контекста, ни способ вызова не влияет
Типовой вывод <<<P000184>,
Но если сам метод расширения является общим,
тип контекста может повлиять на вызов члена.%
?


<<<P001209> Специфика расширения.
<<<P001036>

<<P000978>%
<<<P001057>%
Когда <<<P001210> E}{1}{k} <<<P000185>, являются расширениями
которые доступны и применимы к члену вызова <<<P000186>
(<<P000697>),
Мы определяем понятие
<<<P001211>{специфичность}{расширение!специфичность}
который является частичным порядком на \List{E}{1}{k}.

<<<P001213>{%>
Специфика используется для определения того, какой метод расширения выполнить.
в ситуации, когда возможен более одного выбора.%
?

<<P000979>%
<<<P001058>%
<<<P000187> быть участником вызова с получателем <<<P000188>
и соответствующее имя члена <<<P000189>,
<<<P001059>%
и пусть <<<P000190> и <<<P000191> обозначает два различных
Доступные и применимые расширения для <<<P000192>.
<<<P001060>%
Пусть <<<P000193> быть инстанцирован <<<P001214> Тип <<<P000194> в отношении <<<P000195>,
<<<P000196> быть привязанным к инстанциации <<<P001215> Тип <<<P000197> в отношении <<<P000198>,
<<<P000199>
(P000698).
<<<P000200> Более конкретно, чем <<P000201> В отношении <<P000202>
если выполнено хотя бы одно из следующих условий:

<<P000644>
<<<P001216>
<<P000203> не указывается в системной библиотеке,
<<<P000204> Объявляется в системной библиотеке.
<<<P001217>
<<<P000205> и <<P000206> Обоих заносят в системную библиотеку,
или ни один из них не объявлен в системной библиотеке, и

<<P000645>
<<<P001218> \SubtypeNE{T_1}{T_2}, но не <<P001220>>>{T_2}{T_1}, или
<<P001221> \SubtypeNE{T_1}{T_2} \SubtypeNE{T_2}{T_1} и
\SubtypeNE{S_1}{S_2}, но не \SubtypeNE{S_2}{S_1}.
<<P000662>

<<<P001226>{%>
Другими словами, инстанциированный \ON{} тип определяет специфику,
и связанное с инстанциацией <<<P001228> Тип используется в качестве галстука
в случае, когда субтипирование не проводит различия между первым.%
?
<<P000663>

<<<P001229>{%>
Следующие примеры иллюстрируют неявную резолюцию о продлении
когда доступно несколько расширений.%
?

<<P000003>

<<<<P001230><%
Здесь применяются оба расширения,
<<<P000802> Более конкретно, чем <<<P000803> Потому что
\SubtypeNE{<<P000804>>>}{<<P000805>>>.%
?

<<P000004>

<<<P001232>{%>
Здесь все три расширения относятся к обоим вызовам.
Для <<P000806>, <<P000807> Наиболее специфический,
<<<P000808>> Является подтипом обоих
<<P000809> и <<P000810>.
Следовательно, тип <<<P000811> является <<<P000812>.

Для <<P000813>, <<P000814> является наиболее конкретным.
Презентация \ON{} Типы, которые сравниваются
<<<P000815> <<<P000816> и
P000817 для двух других расширений.
Используя привязку к инстанциации \ON{} Типы как галстук,
Мы находим, что <<<P000818> менее точным, чем <<<P000819>,
Поэтому выбирается <<P000820>.Следовательно, тип <<<P000821> является <<<P000822>.%
?

<<<P001235>{%>
Определение специфики направлено на выбор
Расширение, которое имеет более точную информацию о типе.
Это не обязательно приводит к наиболее точному типу результата.
(перенаправлено с «P000823») мог бы вернуться <<<P000824>,
Также важно, чтобы правило было простым.

На практике мы ожидаем, что непреднамеренные конфликты с именами участников расширения будут редкими.
Если автор предоставляет более специализированные версии
расширение для подтипов,
Выбор расширения, которое имеет наиболее точные типы
Скорее всего, это будет довольно неожиданное и полезное поведение.
?


<<<P001236> Статический анализ членов продления
<<P001037>

<<<P000980>%
Статический анализ деклараций членов в расширении <<<P000207>
опирается на объемы расширения
(<<P000699>)
и соблюдает обычные правила, за исключением следующих:

<<<P000981>%
При проведении статического анализа на теле экземпляра члена
расширение <<<P000208> с \ON{} тип <<P000209>,
Статический тип <<<P001238><<<P000210>.

<<P000982>%
Ошибка времени компиляции возникает, если тело члена расширения
содержит <<P001239>.

<<<P001240>{%>
Лексический поиск в расширении <<<P000211> может давать
Заявление метода экземпляра, заявленного в <<<P000212>.
Как указано в другом месте
(<<P000700>),
Это означает, что члены экземпляра расширения
Теневые члены класса
при вызове из другого экземпляра члена внутри того же расширения
Использование вызова неквалифицированной функции
(перенаправлено с «P000825») и не <<P000826>,
<<P000701>.
%
Это единственная ситуация, когда неявные призывы к
член расширения с именем базы <<<P000950>
может быть успешным, даже если интерфейс приемника имеет
Член с базовым именем <<P001241>.
%
С другой стороны, это согласуется с общим свойством Дарта, что
Лексически заключающие декларации теневые другие декларации, например,
Унаследованная декларация может быть скрыта глобальной декларацией.
Вот пример: %
?

<<P000005>

<<<P001242>{%>
<<<P000827>,
<<P000828> принимает решение по декларации в <<<P000829>,
Если учесть, что <<P000830> не имеет члена с базовым именем <<<P000831>,
Если нет других расширений, создающих конфликт.

Использование <<P000832> в декларации <<P000833>
не определяется в текущем лексическом разрезе,
Таким образом, он рассматривается как <<<P000834>,
потому что интерфейс <<<P001243> тип <<P000835> иметь
<<P000836> Геттер.

Использование <<<P000837> в \code{isOdd} Лексически разрешает
<<P000839> выше него,
Таким образом, он рассматривается как <<<P000840>,
Даже если есть другие расширения в области
где <<<P000841> на <<<P000842>.

Использование <<<P000843> в <<<P000844> Лексически разрешает
<<<P000845> выше него,
Несмотря на то, что есть член с тем же именем в
Интерфейс <<<P001244>.
Геттер <<<P000846> нельзя назвать
в неявном призыве из любого места за пределами <<P000847>.
Это исключительный случай, упомянутый выше,
где член расширения тени
Постоянный член экземпляра на <<<P001245>.
На практике продления очень редко вводят членов.
с тем же именем, что и член его <<<P001246>< Тип интерфейса.

Неквалифицированный идентификатор <<<P000848> который не входит в сферу
Называется «P000849». Члены инспекции, как обычно
(<<P000702>).
Если <<P000850> не определяется статичным типом <<<P001247>
(тип <<<P001248>{})
Это может быть заблуждением.
или может быть решена с использованием другого расширения.%
?


<<<P001249> Метод расширения Клозуризация
<<P001038>

<<P000983>%
Метод расширенного экземпляра подлежит клозуризациианалогично методам классовых экземпляров
(<<P000703>).

<<<P000984>%
<<<P001061>%
Ссылка на стр. Ссылка на стр. Заявка должна быть выполнена с помощью приложения
(<<P000704>)
форма \code{$E$<\List{ S}{1}{m}>($e_1$)}.
<<<P001062>%
<<<P001252>> Y}{1}{m} - формальные параметры типа
Расширение <<<P000216>.
%
<<<P001063>%
Выражение <<P000217> формы <<P000851>
где <<P000951> Если идентификатор известен как
<<<P001253>{добыча имущества расширения}{%
Расширение! извлечение имущества.
Ошибка времени компиляции, если только <<P000219> Заявляет о члене экземпляра по имени <<<P001254>.
Если указанный экземпляр является методом, то
<<P000220> имеет статический тип <<<P000221>,
где <<P000222> Функциональный тип указанного способа декларирования.

<<<P001255>{%>
Если <<P000952> является геттером <<<P000223> Это призыв Геттера,
который указан в другом месте
(<<P000705>) %
?

<<<P000985>%
Если <<P000953> Это метод <<<P000224> Он известен как
<\IndexCustom{внутренний метод клозулизации}{%
Расширение! Пример метода клозурии
<<<P000954> на <<<P000225>,
и оценка <<<P000226>
\commentary{(который <<P001258>>>{$E$<\List{) S}{1}{m}>($e_1$). <<<P001260>> ?
идет следующим образом:

<<<P000986>%
Оценить <<<P000229> на объект <<<P000230>.
<<P000231> Будьте свежей конечной переменной, связанной с <<P000232>.
Далее <<P000233> оценивает объект функции, эквивалентный:

<<P000646>
<<P001261>
<<P000013>
где <<P000955> объявляет параметры типа
<<<P001262>,
требуемые параметры <<<P001263>{p}{1}{n}
Названные параметры <<<P001264>{n}{n}{n}{n}{n}{n}{n}}} с дефолтами <<P001265 >>{1}{k}
использовать <<P000852> для параметров, значение по умолчанию которых не указано.
<<P001266>
<<<P000014>
где <<P000956> объявляет параметры типа
<<<P001267>,
требуемые параметры \List{p}{1}{n}
и факультативные позиционные параметры
\List{p}{n+1}{n+k) с дефолтами <<<P001270>{d}{1}{k}
использовать <<P000853> для параметров, значение по умолчанию которых не указано.
<<P000664>

<<<P000987>%
В функции буквы выше,
<<P000234>,
<<P000235>,
где <<P000236> является типом соответствующего параметра в
Декларация <<<P001271>.
Следует избегать захвата переменных типа в <<<P001272>{S}{1}{m},
<<<P000237> должны быть переименованы, если <<<P001273>{S}{1}{m} содержит любые случаи <<<P000238>,
Для всех <<P000239>.

<<<P001274>{%>
Другими словами, клозуляция является ценностью
функция, буквальная подпись которой совпадает с подписью <<<P001275>,
За исключением того, что фактические аргументы заменяются
Формальные параметры типа <<<P000240>,
Затем он просто передает призыв
<<<P000957>> с захваченным объектом <<<P000241> как получатель.%
?

<<<P001276>{%>
Клосуризация двух расширений никогда не бывает равной
Если только они не идентичны.
Заметим, что это отличается от клозулизации классовых методов экземпляра,
которые равны, когда они отрывают один и тот же метод одного и того же приемника.
?

<<<P001277>{%>
Причина этого различия заключается в том, что даже если <<<P000242> и <<<P000243>
являются примерными методами клозулизации того же расширения <<<P000244>
применяется к тому же приемнику <<<P000245>,
могут иметь различные фактические типы аргументов, переданные <<<P000246>,
потому что аргументы этого типа определяются сайтом вызова
(и с выводом: по статическому типу выражения, дающего <<P000247>),
И не только по свойствам <<P000248> и вырванный метод.%
?

<<<P001278>{%>
Обратите внимание, что метод клозулизации экземпляра на расширении не является
Постоянное выражение,
Даже в том случае, когда приемник является постоянным выражением.
Это потому, что он создает новый объект функции каждый раз, когда он оценивается.
?

<<<P000988>%Клозуризация методом расширения может происходить для неявного вызова
Член расширенного экземпляра.

<<<P001279>{%>
Это является следствием того факта, что неявное обращение
рассматривается как соответствующее явное обращение
(<<P000706>).
Например,
<<P000854> могут быть косвенно преобразованы в
<<P000855>,
которые затем обрабатываются, как указано выше.%
?

<<P000989>%
Клозуризация методом расширения подвергается инстанцированию общей функции
(<<P000707>).
<<<p001280>> Например: }

<<P000006>

<<<P001281>{%>
В этом примере \code{\{\}> напечатаны,
потому что функция объекта получена методом расширения клозуризации
подвергался инстанциации общей функции
В результате получился <<P000856> значение <<<P000857>,
что делает оценку «<<P000858>>>» ложной.%
?


<<<P001283> <<<P001284> Участник продления.
<<P001039>

<<<P000990>%
Расширение может содержать <<<P001285>{} метод, который используется неявно,
Аналогично вызову выражения функции
(<<P000708>).

<<<P001286>{%>
Например, <<<P000859> <<<P000860>
когда статический тип <<<P000255> Это нефункция, которая имеет метод под названием <<<P001287>.
Вот пример, где <<<P001288> Метод происходит от расширения: %
?

<<P000007>

<<<<<p001289>>
Это может показаться несколько удивительным,
Похожий подход с использованием <<<P000861>:
\code{<<P001291>>>{} (<<<P001292 >>) i <<<<P001293> 1[3])\,\,\{\,...\,\}}.
Мы ожидаем, что разработчики будут использовать эту мощность ответственно.
?

<<<P000991>%
<<<P001064>%
<<<P000256> быть приложением расширения
(<<P000709>),
<<P000257> Выражение формы

<<<P001294>
<<P000862>

<<P001295>
<<<P001296>{(где список аргументов типа опущен, когда <<P000263> равен нулю)}.
<<P000264> Затем рассматривается как
(<<P000710>)

<<P001297>
<<P000863>.

<<<P001298>{%>
Другими словами, вызов заявки на продление
немедленно рассматривается как вызов метода расширения, названного <<<P001299>>%
?

<<P000992>%
<<<P001065>%
<<<P000270> выражение со статическим типом <<<P000271>
который не является выражением извлечения свойств
(<<P000711>),
<<<P000272>> быть выражением формы

<<P001300>
<<P000864>

<<P001301>
<<<P001302>{(где список аргументов типа снова опущен, когда <<P000278> равен нулю)}.
Если <<<P000279> <<<P001303>, <<<P001304>, или тип функции,
или интерфейс <<<P000280> имеет метод под названием <<<P001305>,
<<P000281> указывается в другом месте
(P000712).
В противном случае, если <<<P000282> имеет неметодический экземпляр с базовым именем \CALL{}
<<<P000283> - ошибка времени компиляции.
В противном случае <<P000284> Выражение <<P000285> который является

<<P001307>
<<P000865>.

<<<P001308>{%>
Обратите внимание, что <<<P000291> может быть неявным призывом к
метод расширения, названный <<<P001309>,
И это может быть ошибкой.
В последнем случае сообщения об ошибках должны быть сформулированы в терминах <<<P000292>,
не <<P000293>%
?

<<P000993>%
Это ошибка времени компиляции, если <<<P000294> является неявным призывом к
метод экземпляра расширения, названный <<<P001310>.

<<<P001311>{%>
В частности, <<<P000295> Не может быть призывом к приемнику продления
Тип возврата - тип функции <<<P001312> или <<<P001313>.%
?

<<<P001314>{%>
Обратите внимание, что нет поддержки для неявного извлечения имущества.
который отрывает метод расширения, названный <<<P001315>.
Например, при условии расширения <<<P000296> в предыдущем примере: %
?

<<P000008>

<<<P001316>{%>
Неявное извлечение собственности может быть разрешено,
Но это будет стоить читабельности.
Тип <<<P000866> Хорошо известен как не вызывающий,и неявный <<P000867> Отрыв не будет иметь видимого синтаксиса.
Аргументы видны читателю,
Но для неявного отрыва от <<<P001318> функция,
Не существует видимого синтаксиса.%
?

<<<P001319>{%>
При желании извлечение имущества может быть выражено явно.
<<<P000868>>%
?


<<<P001320> Подсчеты
<<P001040>

<<<P000994>%
<<<P000748>, или <<P000749>, используется для представления
фиксированное число постоянных значений.

<<P000647>
<enumType> ::= \ENUM{} <идентификатор>
<<<P001322 >>{} <enumEntry> (',' <enumEntry>)* (',')? "

<enumEntry> ::= <metadata> <identifier>
<<P000665>

<<<P000995>%
Декларация перечня форм
<<<P001323><<<P000297> \ENUM{} <<P000298> \{$m_0\,\,\id_0, \ldots,\ m_{n-1}\,\,\id_{n-1}$\}
имеет тот же эффект, что и классификация

<<<P000015>

<<<P001325>{%>
Это также ошибка времени компиляции для подкласса, смешивания или реализации перечня.
или в явном виде вводить в заблуждение.
Эти ограничения даются в нормативной форме.
в разделах <<<P000713>, <<P000714>, <<P000715>
<<P000716> соответственно.%
?


<<<P001326> Дженерикс
<<P001041>

<<P000996>%
Декларация класса (<<P000717>),
микширование (<<P000718>),
расширение (<<P000719>),
псевдоним (<<<P000720>),
или функция (<<P000721>) <<P000300> может быть <<<P000750>,
То есть <<<P000301> Могут быть объявлены формальные параметры типа.

<<<P000997>%
Когда объект в этой спецификации описывается как общий,
и рассматривается особый случай, когда число аргументов типа равно нулю,
Список аргументов типа должен быть опущен.

<<<P001327>{%>
Это позволяет включать необычные случаи в качестве особых случаев.
Например,
вызов негенерической функции возникает как частный случай
где функция принимает аргументы нулевого типа,
Приводятся аргументы нулевого типа.
В этой ситуации некоторые операции также опускаются (не имеют эффекта).
операции, в которых формальные параметры типа заменяются
Типовые аргументы.%
?

<<<P000998>%
A <<<P001328>{общая декларация класса}{объявление класса}! общий
вводит общий класс в библиотечный объем текущей библиотеки.
A <<<P001329>{общий класс}{класс!общий класс}
Это отображение, которое принимает список фактических аргументов типа.
и относит их к одному классу.
Рассмотрим типовую декларацию класса <<P000302> Название: P000303
с официальными декларациями параметров типа
<<P000304>,
и параметризованный тип <<P000305> формы <<<P000869>.

<<P000999>%
Ошибка времени компиляции <<<P000308>.
Ошибка времени компиляции <<<P000309> не является хорошо ограниченным
(<<P000722>).

<<<P001000>%
В противном случае указанный параметризованный тип <<P000870> обозначает
Приложение общего класса, объявленное <<<P000312> Типовые аргументы
<<P000313>.
Это дает класс <<<P000314> члены которого эквивалентны членам
Классовая декларация, полученная из декларации <<P000315> путем замены
каждое появление <<<P000316> по <<<P000317>.

<<<P001330>{%>
Другие свойства <<P000318> Например, отношения подтипа
указывается в других
(<<<P000723>)%
?

<<<P001331>{%>
Общие псевдонимы типа указаны в другом месте
(<<P000724>>>) %
?

<<<P001001>%
A <<<P001332>{общий} Тип {type!generic} - тип, который вводится
декларация общего класса или псевдоним общего типа,
Или это тип <<P000871>.

<<<P001002>%
A <<<P001333>{объявление общей функции}{объявление функциональности!объявление общей функции}
Вводится общая функция (<<P000725>) в существующий объем.

<<<P001003>%
Рассмотрим функцию вызова выражения формы
<<P000872>,
где статический тип <<<P000873> Является общим типом функции
Формальные параметры типа
<<P000320>.
Ошибка времени компиляции <<<P000321>.
Ошибка времени компиляции при наличии <<<P000322>Например, <<P000323> Не является подтипом <<<P000324>.

<<<P001334>{%>
То есть, если количество аргументов неправильное,
или если аргумент <<P000325> фактического типа не является подтипом
соответствующие границы,
где каждый формальный параметр типа заменяется
Соответствующий аргумент фактического типа.%
?

<<P000648>
<typeParameter> ::= <metadata> <identifier> (<<<P001335>{} <typeNotVoid>)?

<typeParameters>::= <<typeParameter> (>typeParameter>) "
<<P000666>

<<<P001004>%
Параметр типа <<<P000326> может быть суффиксирован с <<<P001336>{} пункт
Это означает, что <<<P000751> для <<<P000327>.
Если нет <<<P001337> Предложение присутствует, верхняя граница <<<P000874>.
Это ошибка времени компиляции, если параметр типа является супертипом его верхней границы.
Когда верхняя граница сама по себе является переменной типа.

<<<P001338>{%>
Это не позволяет делать такие заявления, как
<<<<P001339>< \EXTENDS{} X
и
<<<P001341>{X} <<<P001342>< Y, Y <<<P001343> X}%
?

<<<P001005>%
Параметры типа объявлены в области параметров типа класса или функции.
Типовые параметры дженерика <<<P000328> являются по объему
- границы всех типовых параметров <<<P000329>.
Типовые параметры типовой декларации классов <<<P000330> Кроме того, в сфере применения
\EXTENDS{} и <<P001345>>>{} Положения <<<P000331> (если они существуют)
и в теле <<P000332>.

<<<P001346>{%>
Однако параметр типа общего класса
Считается неполноценным типом
когда на него ссылается статический член
(P000726).
Объемы, связанные с параметрами типа общей функции
Описаны в (<<<P000727>.%)
?

<<<P001347>{%>
Ограничение статических элементов необходимо, поскольку
переменная типа не имеет значения в контексте статического элемента;
Потому что статика делится между
Все генерические инстанциации генерического класса.
Однако переменная типа может ссылаться на инициализатор экземпляра,
Даже если <<<P001348> Недоступно.%
?

<<<P001349>{%>
Поскольку параметры типа находятся в пределах их объема,
Мы поддерживаем F-связанную количественную оценку.
Это позволяет код проверки типа, такой как:%
?

<<P000009>

<<<P001350>{%>
Даже в тех случаях, когда параметры типа находятся в пределах
В настоящее время существует множество ограничений:

<<P000649>
<<<P001351>[<<P000333>]
Параметр типа не может использоваться для обозначения конструктора в
выражение создания экземпляра (<<P000728>).
<<<P001352>[<<P000334>>>]
Параметр типа не может использоваться в качестве суперкласса или суперинтерфейса.
(<<P000729>, <<P000730>, <<P000731>.
<<<P001353>[<<P000335>>>]
Параметр типа не может использоваться в качестве общего типа.
<<P000667>

Приведены нормативные версии этих
в соответствующих разделах данной спецификации.
Некоторые из этих ограничений могут быть сняты в будущем.
?


<<<P001354> Разнообразие
<<P001042>

<<<P001006>%
Мы говорим, что тип <<<P000336> <<<P000752> в виде <<P000337> если {f}
<<P000338> происходит в ковариантном положении в <<<P000339>,
но не в противоположном положении,
И не в инвариантном положении.

<<<P001007>%
Мы говорим, что тип <<<P000340> <<P000753> в виде <<<P000341> если {f}
<<P000342> происходит в противоположном положении в <<<P000343>,
но не в ковариантном положении,
И не в инвариантном положении.

<<<P001008>%
Мы говорим, что тип <<<P000344> <<P000754> в виде <<P000345> если {f}
<<P000346> происходит в инвариантном положении в <<<P000347>,
<<P000348> происходит в ковариантном положении, а также в противоположном положении.

<<<P001009>%
Мы говорим, что тип <<<P000349> встречается в <<<P000755> в виде <<P000350>
Если верно одно из следующих условий:

<<P000650>
<<P001355> <<<P000351> <<<P000352>

<<P001356> <<P000353> Форма: <<P000875>
где <<P000356> обозначает общий класс<<P000357> происходит в ковариантном положении в <<<P000358> Для некоторых <<P000359>.

<<P001357> <<P000360> имеет форму
<<P000876>
где список параметров типа может быть опущен,
<<P000364> происходит в ковариантном положении в <<<P000365>.

<<P001358> <<P000366> имеет форму

<<P000877>

\code{<<P001360>>>($S_1<<P001550>>>x_1, <<P001361>>>\  S_k\ x_k
<<P000369> S_{k+1}\ x_{k+1} = d_{k+1}, \ldots\  S_n\ x_n = d_n$]]

<<P001363>
или в форме

<<P000878>

\code{<<P001365>>>($S_1<<P001556>>>x_1, <<P001366>>>\  S_k\ x_k
<<P000372> S_{k+1}\ x_{k+1} = d_{k+1}, \ldots\  S_n\ x_n = d_n$\}]

<<P001368>
где список параметров типа и каждое значение по умолчанию могут быть опущены,
<<P000373> происходит в противоположном положении в <<<P000374>
Для некоторых <<P000375>.

<<P001369> <<P000376> Форма: <<P000879>
где <<P000379> обозначает общий тип псевдонимов, такой, что
<<P000380>,
Формальный параметр типа, соответствующий <<<P000381> является ковариантным,
<<P000382> происходит в ковариантном положении в <<<P000383>.

<<P001370> <<P000384> Форма: <<P000880>
где <<P000387> обозначает общий тип псевдонимов, такой, что
<<P000388>,
Формальный параметр типа, соответствующий <<<P000389> является противоположным,
<<P000390> Происходит в противоположном положении в <<<P000391>.
<<P000668>

<<<p001010>%
Мы говорим, что тип <<<P000392> встречается в <<<P000756> в виде <<P000393>
Если верно одно из следующих условий:

<<P000651>
<<P001371> <<P000394> Форма: <<P000881>
где <<P000397> обозначает общий класс
<<P000398> происходит в противоположном положении в <<<P000399>
Для некоторых <<P000400>.

<<P001372> <<P000401> имеет форму
<<P000882>
где список параметров типа может быть опущен,
<<P000405> происходит в противоположном положении в <<<P000406>.

<<P001373> <<P000407> имеет форму

<<P000883>

\code{<<P001375>>>($S_1<<P001563>>>x_1, <<P001376>>>, <<P001564>>>S_k<<P001565>>>x_k,
<<P000410> S_{k+1}\ x_{k+1} = d_{k+1}, \ldots\  S_n\ x_n = d_n$]]

<<P001378>
или в форме

<<P000884>

\code{<<P001380>>>($S_1<<P001569>>>x_1, <<P001381>>>\  S_k\ x_k
<<P000413> S_{k+1}\ x_{k+1} = d_{k+1}, \ldots\  S_n\ x_n = d_n$\}]

<<P001383>
где список параметров типа и каждое значение по умолчанию могут быть опущены,
<<<P000414> происходит в ковариантном положении в <<<P000415>
Для некоторых <<P000416>.

<<P001384> <<P000417> Форма: <<P000885>
где <<P000420> обозначает общий тип псевдонимов, такой, что
<<P000421>,
Формальный параметр типа, соответствующий <<<P000422> является ковариантным,
<<<P000423> происходит в противоположном положении в <<<P000424>.

<<P001385> <<P000425> Форма: <<P000886>
где <<P000428> обозначает общий тип псевдонимов, такой, что
<<P000429>,
Формальный параметр типа, соответствующий <<<P000430> является противоположным,
<<P000431> происходит в ковариантном положении в <<<P000432>.
<<P000669>

<<<P001011>%
Мы говорим, что тип <<<P000433> встречается в <<<P000757> в виде <<<P000434>
Если верно одно из следующих условий:

<<P000652>
<<P001386> <<P000435> Форма: <<P000887>
где <<P000438> обозначает общий класс или общий тип псевдонима,
<<P000439> происходит в инвариантном положении в <<<P000440> Для некоторых <<P000441>.

<<P001387> <<P000442> имеет форму

<<P000888>

\code{<<P001389>>>($S_1<<P001576>>>x_1, <<P001390>>>\  S_k\ x_k<<P000445> S_{k+1}\ x_{k+1} = d_{k+1}, \ldots\  S_n\ x_n = d_n$]]

<<P001392>
или в форме

<<P000889>

\code{<<P001394>>>($S_1<<P001582>>>x_1, <<P001395>>>\  S_k\ x_k
<<P000448> S_{k+1}\ x_{k+1} = d_{k+1}, \ldots\  S_n\ x_n = d_n$\}]

<<P001397>
где список параметров типа и каждое значение по умолчанию могут быть опущены,
<<P000449> происходит в инвариантном положении в <<<P000450>
для некоторых <<P000451>,
<<<P000452> Происходит в <<<P000453>
Для некоторых <<P000454>.

<<P001398> <<P000455> Форма: <<P000890>
где <<P000458> обозначает общий тип псевдонима,
<<P000459>,
Формальный параметр типа, соответствующий <<<P000460> является инвариантным,
<<P000461> Происходит в <<<P000462>.
<<P000670>

<<<P001012>%
Рассмотрим заявление о псевдонимах общего типа <<P000463>
с официальными декларациями параметров типа
<<P000464>,
с правой стороны <<P000465>.
Пусть <<<P000466>.
%
Мы говорим, что
<<P001399> Формальный параметр типа <<P000467> является инвариантным {%}
Типовой параметр! инвариантный
f <<<P000468> происходит инвариантно в <<<P000469>, <<P000470>
<<<P001400>{является ковариантным}{параметр типа! ковариантный
f <<<P000471> происходит ковариантно в <<<P000472> и <<<P000473>
\IndexCustom{является контравариантным}{параметр типа!контравариантный}
f <<<P000474> Происходит противоположно в <<<P000475>.

<<<P001402>{%>
Разница дает характеристику того, как тип варьируется
как значение субтермина варьируется, например, переменная типа:
Предположим, что <<P000476> Тип, в котором переменная типа <<<P000477> происходит,
<<<P000478> и <<<P000479> Такие типы, как <<<P000480> Это подтип <<<P000481>.
Если <<P000482> происходит ковариантно в <<<P000483>
<<<P000484> Это подтип <<<P000485>.
Точно так же, если <<<P000486> происходит противоположно в <<<P000487>
<<<P000488> Это подтип <<<P000489>.
Если <<P000490> происходит неизменно
<<<P000491> и <<P000492> не гарантируется
быть подтипами друг друга в любом направлении.
Короче говоря: с ковариацией тип ковариаций;
с противопоставлением, типом противопоказаний;
с инвариантностью, все ставки выключены.%
?


<<<P001403> Суперсвязанные типы
<<P001043>

<<<P001013>%
В этом разделе описывается, как
Заявленные верхние границы формальных типовых параметров соблюдаются.
включая случаи, когда допускается ограниченная форма нарушения.

<<<P001014>%
<<<P000758> Тип <<<P000493> <<<P000891>> Это подтип <<<P000494>.
<<<P001404>{%>
Например, <<<P000892>, <<<P001405> и \VOID{} являются топовыми типами,
<<<P000893> и <<<P001407> Будущее или будущее Или <\DYNAMIC>{}>.%
?

% Мы определяем свойство быть регулярными для всех типов,
% сверхсвязанные для параметризованных типов и хорошо ограниченные
% для всех типов. Мы требуем, чтобы все типы были хорошо ограничены.
% охватывает каждый субтермин типа, который сам по себе является типом, и тогда мы
% требует, чтобы типы были регулярными при использовании в определенных
% ситуаций.

<<<P001015>%
Каждый тип, который не является параметризованным, является <<P000759>.

<<<P001409>{%>
В частности, каждый необычный класс и каждый тип функций.
Регулярно ограниченный тип.%
?

<<<P001016>%
<<<P000495> быть параметризованным типом формы
<<P000894>
где <<P000498> обозначает общий класс или псевдоним общего типа.
Давай.
<<P000895>
Формальные декларации параметров типа <<<P000500>.
<<<P000501> <<<P000760> если {f}
<<P000502> является подтипом
<<P000503>,
Для всех <<P000504>.

<<<<P001410><%
Это означает, что обычные типы являются такими типами.
которые не нарушают предельные значения параметров типа.%
?

<<<P001017>%
<<<P000505> быть параметризованным типом формы
<<P000896>где <<P000508> обозначает общий класс или псевдоним общего типа.
<<<P000509> <<<P000761> Если оба условия истинны:

<<P000653>
<<P001411>
<<<P000510> не имеет регулярных ограничений.
<<<P001412>
<<<P000511> быть результатом замены каждого события в <<<P000512>
верхнего типа в ковариантном положении по <<<P000897>,
и каждое происшествие в <<P000513>
<<<P000898> в противоположном положении по <<<P000899>.
Затем требуется, чтобы <<<P000514> был регулярным.
%
Более того, если <<<P000515> обозначает общий тип псевдонима с телом <<P000516>,
Необходимо, чтобы каждый тип, который возникает как субтермин
<<P000517>
Он хорошо сложен (определяется ниже).
<<P000671>

<<<P001413>{%>
Короче говоря, по крайней мере, один тип аргумента нарушает свои границы, но тип
регулярно ограниченный после замены всех случаев экстремального типа
противоположный экстремальный тип, в зависимости от их дисперсии.%
?

<<<P001018>%
Тип <<<P000518> <<<P000762> если {f}
Он либо регулярный, либо супер-ограниченный.

<<<P001019>%
Любое использование типа <<P000519> Что не очень хорошо ограничено, так это ошибка времени компиляции.

<<<P001020>%
Ошибка времени компиляции при параметризированном типе <<<P000520> сверхограниченный
когда он используется любым из следующих способов:

<<P000654>
<<<P001414> <<P000521> Непосредственная субтерминация нового выражения
(<<P000732>)
Постоянное выражение объекта
(P000733).
<<<P001415> <<P000522> является немедленным субтермином перенаправляющего заводского конструктора
подпись
(P000734).
<<<P001416> <<P000523> является немедленным субтермином <<<P001417> пункт о классе
(<<P000735>),
или это происходит как элемент в списке типов <<<P001418> пункт
% % TODO(Eernst): Приходите с расширением, добавьте ref. Может быть, класс микширования?
(<<P000736>, <<P000737>, <<P000738>,
или <<<P001419> пункт
(<<P000739>, <<P000740>,
или это происходит как элемент в списке типов <<<P001420> пункт о смешивании
(P000741).
<<P000672>

<<<P001421>{%>
Это ошибка <<<P001422>>, если возникает сверхограниченный тип.
как немедленный субтермин <<<P001423> пункт
который определяет границу переменной типа
(<<P000742>)%
?

<<<P001424>{%>
Типы членов из сверхсвязанных классов вычисляются с использованием одного и того же
Правила как типы членов от других типов. Типы функциональных приложений
Сверхограниченные типы вычисляются по тем же правилам, что и типы
Функциональные приложения, включающие другие типы. Вот пример: %
?

<<<p000010>

<<<P001425>{%>
При этом <<<P000900> имеет статический тип <<P000901>,
Даже если верхняя граница переменной типа <<<P000902> составляет <<<P000903>.%
?

<<<P001426>{%>
Сверхограниченные типы позволяют выражать информативные общие супертипы.
некоторые типы, общие супертипы которых
В противном случае было бы гораздо менее информативным.
?

<<<P001427>{%>
Рассмотрим, например, следующий класс: %
?

<<<p000011>

<<<P001428>{%>
Без сверхограниченных типов,
Не существует типа <<P000524> что делает <<P000904> Общий супертип
Все виды формы <<P000905>
(отмечая, что все виды должны быть регулярными)
Когда у нас нет понятия сверхограниченных типов.
Поэтому, если мы хотим, чтобы переменная содержала любой экземпляр типа <<P000906> ''
Эта переменная должна использовать <<P000907> или другой топ-тип
как своего рода аннотация,
что означает, что член, как <<<P000908> Не известно, существует ли
(что мы имеем в виду, говоря, что тип «менее информативный»).
?

<<<<P001429>{%>
Мы могли бы ввести понятие рекурсивных (бесконечных) типов и выразить
наименьшая верхняя граница всех типов формы <<<P000909> как
некоторый синтаксис, значение которого может быть приближено
<<P000910>.
%
Однако мы ожидаем, что любая такая концепция в Дарте
повлечет за собой значительные затраты
О разработчикам и реализацияхс точки зрения дополнительной сложности и тонкости,
Поэтому мы решили этого не делать.
Сверхограниченные типы конечны,
Но они предлагают полезное приближение, контролируемое разработчиком
для таких бесконечных типов.%
?

<<<P001430>{%>
Например,
<<P000911> и
\code{C<C<\VOID>{}>
Это типы, которые разработчик может использовать в качестве аннотации типа.
Этот выбор служит обязательством
конечный уровень раскрытия бесконечного типа,
Это позволяет осуществлять определенный контроль
В точке, где разворачивание заканчивается:
%
Если <<P000912> имеет тип \code{C<C\DYNAMIC>{}>>
<<<P000913> имеет тип <<<<P001435><
<<P000914> не имеет ошибки времени компиляции,
Но если <<P000915> имеет тип \code{C<C\VOID>{}>> Тогда уже
<<<P000916> - ошибка времени компиляции.
Таким образом, разработчики могут получить более или менее строгий режим.
выражения, тип которых выходит за пределы данного конечного развертывания.
?


<<<P001438> Обоснование границы.
<<P001044>

<<<P001021>%
В этом разделе описывается, как вычислять аргументы типа.
которые опущены из вида,
или от призыва к общей функции.

<<<P001439>{%>
% % TODO(Eernst): Когда мы добавим спецификацию типа вывода, мы будем корректировать
% - спецификация i2b, позволяющая принимать начальные значения для
% от фактического типа аргументов (то есть <<P000532>) Дано "от звонящего"
% % для некоторых <<<P000533>; вероятно, «звонящий» всегда будет выводом типа; и
% Все остальные части алгоритма остаются неизменными.
% % Я думаю, что будет более запутанным, чем полезным начать писать это сейчас.
%% и это будет небольшая корректировка, когда мы добавим вывод типа. Так
На данный момент мы просто определяем i2b как автономный алгоритм.
Обратите внимание, что предполагается, что вывод типа уже имел место.
(<<P000743>),
Таким образом, аргументы типа не считаются опущенными, если они выведены.
Это означает, что инстанциация для связывания является механизмом резервного копирования,
который будет использоваться, когда нет информации для вывода.%
?

<<<P001022>%
Рассмотрим ситуацию, когда термин <<P000534> формы \synt
обозначает родовую декларацию типа,
Используется как тип или как выражение в прилагаемой программе.
<<<P001441>{%>
Это означает, что аргументы типа принимаются, но не предоставляются.
?
Мы используем фразу
<<P000763> соответственно <<P000764>
Для определения таких терминов.
В следующем мы упоминаем только сырые виды,
Но все говорили о сырых типах
Применяется к выражениям сырого типа очевидным образом.

<<<P001442>{%>
Например, с заявлением <<P000917>,
Оценка экспрессии сырого типа <<<P000918> В теле уступает
Пример класса <<<P000919> <<<P000920>,
потому что <<<P000921> подлежит инстанциации в связанное.
Обратите внимание, что <<<P000922> не является синтаксически выражением,
Тем не менее, доступ к
<<P000923> пример, подтверждающий <<P000924>
без предварительного согласования,
потому что это может быть значение переменной типа.
?

<<<P001443>{%>
Мы можем однозначно определить сырые типы для обозначения
Результат применения генерического типа
к списку неявно представленных фактических аргументов типа,
Связывание — это механизм, который делает именно это.
Это связано с тем, что Дарт не поддерживает и не будет поддерживать более высокие типы.
Например, значение переменной типа <<<P000535> Это будет тип,
не может быть общим классом <<<P000925> как таковой,
и не может применяться к аргументам типа, например, <<P000926>.%
?

<<<P001444>{%>
В типичном случае, когда встречается только ковариация,
Инстанциация в связанное даст <<<P001445>
Все обычные типы, которые можно выразить.
Это позволяет разработчикам рассматривать сырой тип как тип.
который используется для указания того, что
«Актуальные аргументы типа не имеют значения» <<P000537>.%
?<<<P001446>{%>
Например, принимая декларацию
\code{<<P001448>>>{) C<X расширяет число \{\ldots\},
Включение в строку <<<P000927> доходность <<P000928>,
Это означает, что <<P000929> можно использовать для объявления переменной <<<P000930>
значение которого может быть <<<P000931> для \emph{все возможные} значения <<<P000539>.%
?

<<<P001451>{%>
Рассмотрим ситуацию, когда
общий тип alias обозначает тип функции,
Он имеет один тип параметра, который является противоположным.
Обоснование, связанное с псевдонимом этого типа, затем даст <<<P001452>
всех типов, которые могут быть выражены
Разнообразие таких аргументов.
Это позволяет разработчикам рассматривать такой тип псевдонима, используемый в качестве сырого типа.
Тип функции, который позволяет передавать функцию клиентам.
где не имеет значения, какие ценности
для аргумента типа, который ожидает клиент <<<P000540>.%
?
<<<P001453>{%>
Например, с
\code{\TYPEDEF{} F<X> = \FUNCTION(X)
Включение в строку <<<P000932> доходность <<P000933>,
Это означает, что <<P000934> можно использовать для объявления переменной <<<P000935>
чья ценность будет функцией, которая может быть передана клиентам
a <<<P000936> для \emph{все возможные} значения $T$.%
?


<<<P001458> Вспомогательные концепции для обоснования границ
<<P001045>

<<<P001023>%
Прежде чем мы определим моментализацию для связывания
Необходимо определить два вспомогательных понятия.
<<<P000543> Будьте сырым типом.
Тип <<P000544> затем
<<<P001459>{сырая зависимость на}{сырая зависимость на!тип}
<<P000545> если выполняется одно или несколько из следующих условий:

<<P000655>
<<P001460>
<<P000546> имеет форму \synt{typeName}, $S$ <<<P000548>.
<<<P001462>{%>
Обратите внимание, что этот случай не применим
Если <<P000549> является субтермином термина формы
<<<P001463><<<P000550> <typeArguments>,
То есть,
Если <<P000551> Получает любые аргументы.
Также обратите внимание, что <<<P000552> не может быть переменной типа,
Потому что тогда '<<<P000553> не может удержаться.
См. обсуждение ниже и ссылку на ~<<<P000744>
Подробнее о том, почему это так.%
?
<<P001464>
<<P000555> имеет форму <<<P001465>{<typeName> <typeArguments>}
Один из типов аргументов зависит от <<<P000556>.
<<P001466>
<<P000557> имеет форму \syntax{<typeName> <typeArguments>?}<<<P001591> где
<<<P001468>{typeName} обозначает псевдоним типа <<P000558>,
и тело <<P000559> Это зависит от <<<P000560>.
<<P001469>
<<P000561> имеет форму
<<<P001470><<тип> \FUNCTION{} <typeParameters>? <parameterTypeList>
\syntax{<type>?}\ raw-depends on $T$,
или связанное в <<<P001473>{<typeParameters>?}<<P001593>raw-depends на <<<P000563>,
или тип в <<<P001474>{параметрTypeList} зависит от <<<P000564>.
<<P000673>

<<<P001475>{%>
Метаварианты
(<<P000745>)
как <<P000565> и <<P000566> понимаются как обозначающие типы,
и они считаются равными (как в '<<<P000567> является $T$)
в том же смысле, что и в разделе о правилах подтипа
(P000746).
%
В частности,
Хотя два одинаковых синтаксиса могут обозначать два различных типа,
и две разные части синтаксиса могут обозначать один и тот же тип,
Свойство интереса здесь заключается в том, обозначают ли они один и тот же тип.
Не то, чтобы они были написаны одинаково.

Интуиция, стоящая за ситуацией, когда тип зависит от другого типа
Мы должны вычислить любые недостающие аргументы для последнего.
Чтобы понять, что означает первое.

В правиле о псевдонимах типа <<<P000569> может быть, а может и не быть,
Аргументы типа могут присутствовать или не присутствовать.Однако нет необходимости рассматривать результат замены.
фактические аргументы типа для формальных параметров типа в корпусе <<P000570>
(или даже правильность передачи аргументов этого типа на <<P000571>),
Потому что нужно только проверить
Все виды формы <<<<P001476><<Наименование типа>>
И на них не влияет такая замена.
Другими словами, сырая зависимость — это отношение.
Что легко и дешево вычислить.%
?

<<<P001024>%
<<<P000572> быть общим классом или общим типом псевдонимов
<<P000573> Официальные заявления о параметрах типа
содержащие формальные параметры типа <<<P001477> X}{1}{k} и границы <<<P001478> B}{1}{k}
Для любого <<P000574>
Мы говорим, что формальный параметр типа <<<P000575> <<<P000765>
когда выполняется одно из следующих требований:

<<P000656>
<<P001479> <<<P000576>> Опущено.

<<<P001480> <<P000577> включен, но не содержит ни одного из <<<P001481>{X}{1}{k}.
Если <<P000578> сырой зависит от сырого типа <<P000579>
каждый параметр типа <<<P000580> Должен иметь простую границу.
<<P000674>

<<<P001025>%
Понятие простой связи должно толковаться индуктивно, а не
мондуктивно, т. е. если связан <<P000581> общего класса или
общий псевдоним типа <<<P000582> достигается в ходе расследования, является ли
<<P000583> Это простая граница, ответ - нет.

<<<P001482>{%>
Например, с
\code{<<P001484>>>{) C<X\EXTENDS{}C><<P001594>>>\}
Параметр типа <<<P000937> не имеет простой границы:
Необработанный <<P000938> используется в качестве границы для <<<P000939>,
<<<P000940> должны иметь простые границы,
Но одна из границ <<<P000941> является пределом <<<P000942>,
и эта граница <<<P000943>, так что <<<P000944> должны иметь простые границы:
Это был цикл, поэтому ответ «нет».
<<P000945> Не имеет простых границ.%
?

<<<P001026>%
<<<P000584> быть общим классом или общим псевдонимом.
Мы говорим: <<P000585>
<<<P001486>{имеет простые границы}{тип! Общий, имеет простые границы.
f каждый параметр типа <<<P000586> Имеет простые границы.

<<<P001487>{%>
Теперь мы можем указать, в каком смысле инстанциация в связанное требует
Типы должны быть «достаточно простыми».
Мы накладываем следующее ограничение на все границы параметров типа:
потому что все параметры типа могут быть подвергнуты инстанциации в связанные.%
?

<<<P001027>%
Это ошибка времени компиляции
Если формальный параметр типа связан <<<P000587> Содержит сырой тип <<P000588>,
Если только <<P000589> Имеет простые границы.

<<<P001488>{%>
Поэтому аргументы типа на границах могут быть только опущены.
Если они сами имеют простые границы.
В частности,
\code{\CLASS{) C<X\EXTENDS{}C><<P001596>>>\}
Ошибка времени компиляции,
потому что связан <<<P000946> сырой,
и формальный параметр типа <<<P000947>
что соответствует опущенному типу аргумента
не имеет простой границы.%
?

<<<P001028>%
<<<P000590> быть одним из видов формы <<<<P001492>{название типа}
который обозначает общий класс или псевдоним общего типа
<<<P001493><<<P000591> Он сырой.
Тогда <<<P000592> эквивалентен параметризованному типу, который
результат, полученный путем применения инстанциации к привязке к <<P000593>.
Это ошибка времени компиляции, если инстанциация в связанный не удается.

<<<P001494>{%>
Это правило применимо ко всем случаям необработанных видов,
Например, когда это происходит как аннотация типа переменной или параметра,
в качестве функции возврата,
как тип, который тестируется в тесте типа,
как тип в <<<P001495> Часть.
и т.д.%
?


<<<P001496> Обоснование связанного алгоритма
<<P001046>

<<<P001029>%
Теперь уточним, как
<<P000766>
Алгоритм продолжается.
<<<P000594> Будьте сырым типом.
Пусть <<<P001497>{X}{1}{k} будут формальными параметрами типа в декларации <<<P000595>,
И пусть <<<P001498> B {1} {k} - их границы.Для каждого <<P000596>
<<<P000597> обозначать результат инстанциации связанным на <<<P000598>;
В случае, когда граница <<<P000599> опущена, пусть <<<P000600> будет <<<P001499>.

<<<<P001500><%
Если <<P000601> Для некоторых <<<P000602> сырой (в общем: если сырой зависит от какого-то типа); <<P000603>
Тогда все его (соответственно P000604) опущенные аргументы типа имеют простые границы.
Это ограничивает сложность инстанциации для <<<P000605>,
В частности, это не может быть связано с циклом зависимости.
где нам нужен результат от инстанциации, связанный с <<<P000606>
для того, чтобы вычислить инстанциацию, связанную с <<<P000607>.%
?

<<<P001030>%
Пусть <<<P000608> будет <<<P000609>, для всех <<<P000610>.
<<<P001501>{%>
Это "текущее значение" границы для переменной типа <<<P000611>, на шаге 0;
В целом мы рассмотрим текущий шаг <<<P000612> и используем данные для этого шага.
Например, связанный <<<P000613>, чтобы вычислить данные для шага <<<P000614>, %
?

{%%} Область для определений отношений, используемых во время i2b

\def\Depends{\ensuremath{\rightarrow_m)
\def\TransitivelyDepends{<<P001508>>>{<<P001509>>>^{+}_m}}

<<<P001031>%
Пусть <<<P001510>{} будет отношением между переменными типа
\List{X}{1}{k) такой, что
<<P000615> <<<P000616> Происходит в <<<P000617>.
<<<P001512>{%>
Таким образом, каждая переменная типа связана, то есть зависит от:
Каждая переменная типа в своей связи, которая может включать себя.
?
Пусть <<<P001513>{} будет транзитивным
\commentary{(но не рефлексивно)}
Закрытие <<<P001515>.
Для каждого <<<P000618>, пусть $U_{i,m+1}$, для <<<P000620>,
определяется следующим итеративным процессом, где <<<P000621> обозначает
<<P000948>:

<<P000657>
<<<P001516> [1]
Если существует <<P000624> такой, что
<<P000625>
<<<P001517>{(то есть, если граф зависимостей имеет цикл)}
Пусть \List{M}{1}{p} будут сильно связанными компонентами (SCC)
в отношении <<<P001519>.
<<<P001520>{%>
То есть максимальные подмножества <<<P001521> X}{1}{k)
где каждая пара переменных в каждом подмножестве
соотносятся в обоих направлениях по <<<P001522>;
Обратите внимание, что SCC попарно разъединены;
Кроме того, они однозначно определены вплоть до переупорядочения.
и порядок не имеет значения для этого алгоритма.
?
<<<P000626> быть объединением \List{M}{1}{p}
<<<P001524>{(то есть все переменные, участвующие в цикле зависимости)}.
Пусть <<<P000627>.
Если <<P000628> не принадлежит <<P000629> <<<P000630> <<<P000631>.
В противном случае существует <<<P000632> <<<P000633>;
<<P000634> Получено из <<<P000635>
путем замены \DYNAMIC{} для каждого появления переменной в <<<P000636>
Находится в позиции <<<P000637> который не является контравариантным,
и заменить <<P000949> для каждого появления переменной в <<<P000638>
Находится в противоположном положении в <<<P000639>.

